\documentclass[11pt]{article}

\usepackage[margin=1in]{geometry}
\usepackage{booktabs}
\usepackage{tabularx}
\usepackage{array}
\usepackage{enumitem}
\usepackage[hidelinks]{hyperref}
\usepackage{xcolor}

\definecolor{linkblue}{HTML}{075E88}
\hypersetup{
  colorlinks=true,
  linkcolor=linkblue,
  urlcolor=linkblue,
  citecolor=linkblue
}

\setlength{\parindent}{0pt}
\setlength{\parskip}{0.65em}
\renewcommand{\arraystretch}{1.18}

\title{\textbf{Authentication Autopilot Cascade Study}\\
\large Advisor Project Update: Motivation, Research Questions, Methodology,
and Implementation Status}
\author{Kofi Ampomah\\UA SPECTRAL Lab, The University of Alabama}
\date{\today}

\begin{document}
\maketitle

\begin{abstract}
This project investigates a familiar but underexamined security problem: users
are often asked to make consequential authorization decisions only after they
have completed several routine authentication steps. The study tests whether
the momentum created by signing in, continuing through an identity provider,
and approving an MFA request reduces careful review of a subsequent
AI-assistant permission screen. It also examines whether placing a single
over-broad permission at the beginning or end of an otherwise identical list
changes behavioral exposure to, and later recall of, that permission. The
prototype implements a $2 \times 2$ between-subjects experiment, records
behavioral measures through a university-hosted API, and connects each task record to a follow-up
Qualtrics survey. The planned sample consists of adult university students
recruited through Prolific.
\end{abstract}

\section{Introduction}

Digital identity workflows rarely end with a password. A student connecting a
new tool may choose a sign-in method, move through an identity-provider
handoff, approve an MFA prompt, and then encounter a request that allows an AI
assistant to read data or act on the student's behalf. From the system's
perspective, these screens perform distinct functions. From the user's
perspective, however, they can feel like one continuous sequence whose apparent
goal is simply to ``get through'' the login.

That distinction is central to the study. Authentication establishes or
verifies identity; authorization determines what another application or agent
may do after identity has been established. Completing a legitimate sign-in
and MFA check does not make a later permission request proportionate, safe, or
necessary. Yet people make these decisions with limited time and attention.
After several successful clicks, the next approval may feel less like a new
security judgment and more like the final step of a process already underway.
The term \emph{authentication autopilot} is used here to describe this possible
shift from deliberate review toward habitual continuation.

Prior work has documented usability burdens associated with MFA
\cite{reese2019usability}, limited attention and comprehension during
permission decisions \cite{felt2012android}, and susceptibility to deceptive online requests
within academic communities \cite{diaz2020,sheng2010}. University students are
a substantively relevant population rather than merely a convenient sample.
They routinely encounter institutional SSO and MFA, adopt new educational and
productivity services, and increasingly interact with AI-enabled tools that
request access to files, calendars, and other accounts. The study therefore
places a consequential authorization decision within a workflow that is
recognizable to the intended participants, while ensuring that no real account
or credential is used.

\section{Study Motivation}

\subsection{Security decisions occur in sequences}

Security interfaces are often designed and evaluated as isolated moments, but
people experience them as sequences. A sign-in button leads to a provider
handoff; the handoff leads to MFA; MFA leads to a permission request. Each
successful step reinforces the expectation that continuing is normal and that
the next approval belongs to the same trusted process. This interactional
momentum may be especially consequential at the authorization stage, where the
user is no longer proving identity but granting power to another system.

The study is motivated by the possibility that a technically valid
authentication chain can create misplaced confidence in a later request.
Nothing needs to be visibly broken or overtly deceptive. A user may simply
transfer trust from the earlier identity checks to a permission screen that
deserves an independent judgment. Measuring this subtle carryover requires
observing behavior across the whole sequence rather than presenting a consent
screen in isolation.

\subsection{AI delegation raises the consequences of inattentive approval}

AI assistants increasingly connect to calendars, files, messages, and other
services. These connections can be useful, but they may also grant persistent
or action-oriented authority. The difference between allowing an application
to view a profile and allowing an assistant to make purchases is not merely a
difference in data sensitivity; it is a difference in what the system can do
after access is granted.

For this reason, the planted permission in the experiment concerns access to
saved payment methods and the ability to make purchases on the user's behalf.
It is intentionally inconsistent with the stated scheduling-and-study-file
task and creates a clear risk of financial harm. The manipulation therefore
tests attention to a meaningful boundary violation rather than an arbitrary
change in wording.

\subsection{Why university students are the target population}

University students operate in a particularly dense digital environment. They
regularly move among institutional portals, cloud productivity tools, learning
platforms, third-party applications, SSO interfaces, and MFA prompts. They also
have practical reasons to try new educational tools, free services, discounts,
and productivity applications. The rationale for studying students is
therefore their repeated exposure to these workflows, not an assumption that
students are inherently careless.

Recent research further supports treating students as a distinct security
population whose online practices are shaped by high technology use, uneven
formal security training, and context-specific motivations
\cite{gwenhure2025}. At the same time, AI use has become part of ordinary
academic work. One recent study at a selective U.S. college found that more
than 80\% of surveyed students reported academic use of generative AI within
two years of ChatGPT's release \cite{contractor2025}. That estimate should not
be generalized to every institution, but it demonstrates that AI-related
authorization is no longer a hypothetical concern for student populations.

\subsection{Scientific and practical value}

Scientifically, the project connects research on security fatigue, permission
attention, authentication usability, and emerging AI delegation. Its factorial
design makes it possible to separate the effect of the preceding chain from the
effect of permission position and to test whether the two factors compound one
another.

Practically, the findings may inform how AI-assistant consent interfaces are
ordered, timed, and separated from authentication flows. If later placement or
authentication momentum reduces meaningful review, designers may need clearer
transitions, stronger scope explanations, or additional confirmation for
high-risk delegated actions. The objective is not to remove useful automation,
but to support authorization decisions that remain understandable even when
users arrive at them after several routine approval steps.

\section{Related Work}
\label{sec:related}

The project brings together five bodies of research that have largely
developed in parallel: security fatigue, SSO mental models, MFA usability,
permission attention, and AI-assistant delegation. Taken separately, these
literatures explain important parts of the problem. Taken together, they
suggest that the location of an authorization request within a longer
interaction may be as important as the wording of the request itself.

\subsection{Security Fatigue and Sequential Decisions}
\label{sec:rw-fatigue}

Security fatigue describes the resignation, avoidance, and reduced engagement
that can arise when people are repeatedly asked to make security decisions
\cite{stanton2016fatigue}. This account is consistent with broader theories of
effortful judgment \cite{kahneman2011thinking} and with comparative work
showing that authentication mechanisms impose distinct cognitive and usability
costs \cite{bonneau2012quest}. Experimental evidence also shows that repeated
exposure can produce habituation to security warnings. Eye-tracking,
neuroimaging, and longitudinal studies have found declining attention across
repeated warning encounters \cite{anderson2016memory,vance2018tuning}, while
visual attention attractors can reduce, but not eliminate, this decline
\cite{bravolillo2014harder}.

Most of this literature studies repeated versions of the same warning. The
present study examines a different form of repetition: a heterogeneous
sequence in which sign-in, identity-provider handoff, MFA, and AI
authorization are technically distinct but may be experienced as one
continuous approval process.

\subsection{Single Sign-On: Mental Models and Consent}
\label{sec:rw-sso}

SSO is often the first point at which users encounter delegated identity.
Research has documented substantial conceptual confusion about what
information is shared during web SSO and whether the relying application
receives the user's identity-provider password
\cite{sun2013investigating}. Related work has examined why users accept or
reject SSO \cite{sun2011makes}, while more recent research reports widespread
use alongside incomplete understanding and weak account-security practices
\cite{petrie2024passwords}.

These mental-model problems are consequential because OAuth-based sign-in may
also involve requests for profile information or additional scopes. Large-scale
analyses have identified scope over-request in deployed SSO systems
\cite{dimova2023ssomething} and personal-information leakage during
authenticated browsing \cite{pham2025consented}. The consent interface is
therefore the point at which an abstract delegation protocol becomes a
decision that the user must interpret.

\subsection{Multi-Factor Authentication: Usability and Adoption}
\label{sec:rw-mfa}

MFA provides meaningful protection against account compromise, but its
security benefit does not remove its interaction cost. Comparative work has
evaluated the usability of common second-factor methods
\cite{reese2019usability}, and university studies have documented both
adoption and practical friction \cite{colnago2018horrible}. Student-focused
work further suggests that MFA can create frustration and interfere with
time-sensitive activity \cite{arnold2022emotional}. A systematic review of
user perceptions similarly shows that usability, context, and perceived burden
shape MFA acceptance \cite{das2019evaluating}.

The present study does not test whether MFA itself is desirable or secure. It
asks whether completing an MFA approval contributes to interactional momentum
that carries into the next screen. In other words, MFA is treated as one
component of the preceding chain and not as the source of the over-broad
permission.

\subsection{Permission Attention and Comprehension}
\label{sec:rw-permissions}

The permission literature provides the methodological foundation for treating
a consent screen as an observable decision point. Foundational work found low
attention and comprehension at Android install-time permission screens
\cite{felt2012android}. Later studies of runtime dialogs documented low denial
rates and the influence of contextual trust
\cite{bonne2017exploring}, while subsequent research examined users'
expectations and everyday management of privacy and security settings
\cite{frik2022expectations}.

Emerging research extends these questions to LLM-based applications. CLEAR,
for example, uses contextual analysis to help explain privacy-policy risks in
LLM applications \cite{chen2025clear}. A recent preprint on fine-print
injection reports a human baseline in which 97.4\% of participants consented to
malicious policies presented with low visual salience
\cite{chen2025fineprint}. Because that result comes from a recent preprint and
a different task context, it should not be treated as a direct estimate for
the present experiment. It does, however, reinforce the need to examine where
high-risk terms appear and whether users can recall them after approving.

\subsection{AI-Assistant Delegation and Human Oversight}
\label{sec:rw-agents}

AI assistants introduce a qualitatively different authorization problem
because granted access may support persistent, autonomous, or action-oriented
behavior. Technical work on the Model Context Protocol and related tool
ecosystems has identified security, governance, tool-integrity, and
maintainability risks
\cite{hasan2025mcpfirstglance,errico2025securingmcp,gaire2025sokmcp,bhatt2025etdi}.
These mechanisms can define scopes and tokens, but they cannot by themselves
ensure that a person understands the practical consequences of approval.

The human-factors evidence is newer and less settled. Studies of conversational
agents show that perceived safety, trust, and actual risk do not necessarily
align \cite{tolsdorf2025safety}; disclosure research shows that users must
trade perceived benefit against privacy risk
\cite{zhang2024fairgame}; and emerging work argues for privacy management and
authorization models that preserve meaningful human oversight
\cite{xu2026humanunit,humancentered2026authorization}. Industry reports on
consent phishing further demonstrate that attackers can misuse legitimate
OAuth authorization flows to obtain over-broad access
\cite{obsidian2026consentphishing}. These emerging 2025--2026 sources are
included to establish the current direction of the field, but records marked
for verification in the accompanying bibliography must be checked before
publication.

\subsection{Synthesis and Research Gap}
\label{sec:rw-gap}

The literature provides strong reasons to expect reduced attention under
repetition, persistent confusion about SSO, measurable friction in MFA, and
weak comprehension of permission dialogs. What it does not yet explain well is
how these effects combine when an AI-assistant authorization screen appears at
the end of a realistic authentication sequence. Permission-position research
also rarely crosses this full chain.

The present study addresses that gap by manipulating both the preceding
authentication chain and the position of one over-broad permission. Its
intended contribution is a controlled, user-centered account of AI
authorization that combines in-the-moment behavior with immediate recall and
survey measures. Scroll behavior and viewport entry indicate whether
information was available for inspection; they are not presented as eye
tracking or proof of conscious attention.

\section{Research Questions}

\textbf{RQ1: Mental models.} How do university students understand SSO, MFA,
and AI-assistant authorization, and where do misconceptions occur across these
three layers?

\textbf{RQ2: Permission attention.} Do university students encounter and
recall an over-broad AI-assistant permission, and does its position in the
permission list affect review behavior and recall?

\textbf{RQ3: Cascade effect.} Does a preceding sign-in, provider-handoff, and
MFA chain reduce attention to a subsequent AI-assistant consent screen?

\textbf{RQ4: Access management.} What do university students know and believe
about revoking delegated AI-assistant access and identifying actions performed
by an assistant?

These formulations distinguish objective behavior from self-report. In
particular, the implementation can establish that a permission entered the
browser viewport, but viewport entry alone is not proof that the participant
looked at, comprehended, or consciously noticed it.

\section{Methodology}

\subsection{Participants and Recruitment}

The planned population is adults who are currently enrolled at a college or
university. Participants will be recruited through Prolific using an available
student or education prescreener. If the required built-in prescreener is not
available, Prolific's approved custom-screening procedure will be used.
Qualtrics will confirm that the participant is at least 18 years old and
currently enrolled. The exact sample size will be determined through an
appropriate power analysis before launch.

Recruitment, screening, compensation, consent language, and data-retention
procedures will follow the final IRB-approved protocol. The University of
Alabama is the research institution, but enrollment at the University of
Alabama is not an eligibility requirement.

\subsection{Experimental Design}

The behavioral task uses a $2 \times 2$ between-subjects factorial design. Each
participant is assigned to one of four cells and sees only that cell.

\begin{table}[ht]
\centering
\caption{Experimental factors and levels.}
\begin{tabularx}{\textwidth}{>{\bfseries}p{0.18\textwidth}
                            p{0.28\textwidth}
                            X}
\toprule
Factor & Levels & Manipulation \\
\midrule
Authentication chain &
Fresh; Chained &
Fresh participants proceed directly from the study introduction to the
permission screen. Chained participants first complete a simulated sign-in
choice, provider handoff, and number-matching MFA approval. \\

Permission placement &
Top; Bottom &
The over-broad payment permission is rendered either first or last. The six
permissions and their wording are otherwise identical. \\
\bottomrule
\end{tabularx}
\end{table}

\begin{table}[ht]
\centering
\caption{The four experimental cells.}
\begin{tabularx}{\textwidth}{>{\bfseries}p{0.22\textwidth} X X}
\toprule
Cell & Participant experience & Primary comparison \\
\midrule
Fresh + Top &
Permission screen shown without a preceding chain; payment permission first. &
Reference condition for clear early exposure. \\
Fresh + Bottom &
Permission screen shown without a preceding chain; payment permission last. &
Placement effect without the chain. \\
Chained + Top &
Permission screen follows sign-in, provider handoff, and MFA; payment
permission first. &
Chain effect with early placement. \\
Chained + Bottom &
Permission screen follows the complete chain; payment permission last. &
Combined condition used to estimate the interaction. \\
\bottomrule
\end{tabularx}
\end{table}

The principal factorial test is whether the effect of bottom placement becomes
larger after the authentication chain. Directionally, the chained-plus-bottom
cell is expected to produce the lowest exposure and recall, while
fresh-plus-top is expected to produce the clearest opportunity for review.
These are hypotheses, not labels that will be shown to participants.

\subsection{Behavioral Procedure}

\begin{enumerate}[leftmargin=1.6em]
  \item \textbf{Study information and agreement.} The participant reads the
        study information and agrees to begin. The page instructs the participant
        to use only study-provided details and not to enter personal credentials.
        To reduce demand characteristics, the simulated nature of the screens is
        not disclosed until the post-task debrief, subject to final IRB approval.

  \item \textbf{Authentication chain.} In the chained condition only, the
        participant chooses a simulated sign-in method, completes a provider
        handoff, and responds to a simulated phone notification using
        two-digit number matching. The fresh condition skips these screens.

  \item \textbf{AI-assistant permission decision.} TaskFlow AI Assistant asks
        for six permissions. Five are plausible baseline or moderate-risk
        permissions. The planted over-broad permission asks to access saved
        payment methods and make purchases on the user's behalf. The
        participant selects Allow or Cancel.

  \item \textbf{Immediate recall.} The participant selects all permissions
        remembered from a shuffled list containing the six real permissions
        and two foils. The task also asks for the assistant's name and whether
        any requested access seemed broader than necessary.

  \item \textbf{Debrief and survey handoff.} The participant is told that the
        screens were simulated, that no real access was granted, why disclosure
        was delayed, and which permission was intentionally over-broad. After
        the task record is saved, the
        participant continues to Qualtrics. The task passes the anonymous
        participant code, chain condition, and placement condition as
        \texttt{pid}, \texttt{condition}, and \texttt{placement}.
\end{enumerate}

\subsection{Behavioral Measures}

\begin{table}[ht]
\centering
\caption{Behavioral outcomes recorded by the web task.}
\begin{tabularx}{\textwidth}{>{\bfseries}p{0.25\textwidth} X}
\toprule
Measure & Operationalization \\
\midrule
Permission exposure &
Whether the planted payment row was initially visible and whether at least
75\% of the row later entered the permission-list viewport. This is a viewport
exposure measure, not eye tracking. \\

Review behavior &
Whether the permission list was scrolled, maximum scroll percentage, time at
which the planted row entered the viewport, and total decision latency. \\

Authorization decision &
Allow or Cancel, with decision time and combined indicators describing whether
the participant scrolled or was exposed to the planted item before deciding. \\

Permission recall &
Per-item recall for all six real permissions, recall of the planted payment
permission, number of correct items, and false alarms from the two foils. \\

Chain behavior &
Selected sign-in route, screen sequence, time on each screen, back-navigation,
MFA selections, and whether the selected number matched the phone
notification. \\

Secondary responses &
Assistant-name recognition and perceived breadth of the requested access. \\

Completion quality &
Completed versus abandoned status, finish timestamp, and task-to-survey
matching by anonymous participant code. \\
\bottomrule
\end{tabularx}
\end{table}

Empty MFA fields are expected for fresh-condition participants because those
participants never receive an MFA prompt. Empty recall fields are expected in
an abandonment record when the participant leaves before submitting recall.

\subsection{Follow-up Survey}

The Qualtrics instrument currently contains 36 items across eight blocks:

\begin{enumerate}[leftmargin=1.6em]
  \item entry and eligibility;
  \item immediate task reflection;
  \item familiarity with SSO, MFA, and AI assistants;
  \item baseline security behavior;
  \item knowledge and misconceptions across SSO, MFA, and AI authorization;
  \item perceived attention and fatigue across repeated authentication steps;
  \item revocation and auditability; and
  \item demographics and optional closing feedback.
\end{enumerate}

RQ1 is addressed mainly by keyed and open-text survey items concerning SSO,
MFA, permission persistence, and assistant autonomy. RQ2 combines behavioral
exposure and recall with self-reported reading and scrolling. RQ3 is addressed
primarily by the randomized fresh-versus-chained comparison, with survey items
providing contextual self-report. RQ4 is addressed by revocation knowledge,
prior revocation experience, and perceived ability to identify
assistant-initiated actions.

The current auditability item measures perceived ability rather than
demonstrated ability. Consequently, the paper should use language such as
``perceived auditability'' unless an objective scenario-based knowledge item is
added. Open-text responses will be coded using a documented codebook by two
independent coders, with inter-rater agreement reported.

\subsection{Data Architecture and Privacy}

The planned data path is:
\[
\text{Prolific} \rightarrow
\text{university-hosted behavioral task and API} \rightarrow
\text{university-managed database} \rightarrow
\text{Qualtrics} \rightarrow
\text{Prolific completion}.
\]

The university-managed database will contain two tables. \texttt{task\_events} stores completed and
best-effort abandonment payloads. \texttt{task\_withdrawals} stores requests
for data removal. Each row has an automatically generated UUID primary key.
The browser submits validated POST requests to a narrow study API; it has no
database credentials and no public route for selecting, updating, or deleting
records. Hosting, access control, backup, retention, and deletion procedures
will be finalized with the university service owner and documented in the IRB.

The task records a randomly generated participant code, assigned condition,
interaction events, timestamps, browser user-agent string, and general display
dimensions. It does not request or transmit a real password, account, email
address, or financial information. Any connection between a Prolific ID and
the internal participant code must be described in the IRB protocol and
minimized to what is necessary for compensation and record matching.

Abandonment capture is best effort because browsers do not guarantee that a
network request will finish when a page is closed. Completed task submissions
wait for a database response and are therefore more reliable.

\subsection{Planned Analysis}

The primary factorial models will include chain, placement, and their
interaction. Binary outcomes---including Allow versus Cancel, planted-item
viewport exposure, and planted-item recall---can be analyzed using logistic
regression. Continuous or bounded outcomes---including decision latency and
scroll depth---will be examined with distributions appropriate to their scale;
latency may require transformation or a robust model.

The chain-by-placement interaction is central because it tests whether bottom
placement is especially harmful after the authentication sequence. Main
effects will estimate the average effects of chain and placement across the
other factor. Analyses of survey knowledge, perceived fatigue, security
behavior, and prior experience will be secondary or explanatory unless
otherwise preregistered.

Planted-payment recall must not be used as an exclusion criterion because it is
a primary RQ2 outcome. Assistant-name recognition may also be influenced by the
experimental condition and will be treated as a secondary attention or memory
measure rather than a neutral exclusion check. Independent quality rules will
be defined before data collection and may include:

\begin{itemize}[leftmargin=1.6em]
  \item failure to meet age or student eligibility requirements;
  \item incomplete task or missing matching survey record;
  \item duplicate participation under a preregistered rule;
  \item technically invalid records; and
  \item failure of the Qualtrics instructed-response checks, applied according
        to the preregistered rule and Prolific policy.
\end{itemize}

\section{Current Implementation Status}

\begin{table}[ht]
\centering
\caption{Status as of this project update.}
\begin{tabularx}{\textwidth}{>{\bfseries}p{0.29\textwidth} X}
\toprule
Component & Status \\
\midrule
Behavioral interface &
Implemented in a single responsive HTML file with Fresh and Chained routing,
top and bottom placement, simulated sign-in, provider handoff, number-matching
MFA, permission decision, recall, debrief, and withdrawal control. \\

Experimental assignment &
Random assignment across all four cells is active. A restricted
\texttt{force=1} query mode supports controlled testing of each cell. \\

Behavioral instrumentation &
Implemented for screen timing, click sequence, MFA attempts, scroll behavior,
viewport exposure, decision, recall, false alarms, completion, and
abandonment. \\

Database &
The browser client is refactored for a university-hosted, POST-only study API.
Provisioning the approved application host, API, and managed database remains
pending; the earlier Supabase prototype is not the planned participant-data
deployment. \\

Survey &
The 36-item Qualtrics instrument is drafted and mapped to the four research
questions. The eligibility wording now targets college and university
students generally. \\

Completed testing &
Saved payloads have been inspected for fresh and chained conditions, top and
bottom placement, successful MFA matching, completed responses, attention
responses, scrolling, and intentional abandonment. \\

Remaining configuration &
Publish the final Qualtrics survey, replace the placeholder Qualtrics URL,
configure Prolific ID pass-through and completion routing, and deploy the task
to its final HTTPS host. \\

Remaining validation &
Document a complete Fresh + Top run, test every cell end to end through
Qualtrics, perform desktop/mobile/keyboard checks, verify database access
restrictions live, pilot timing, and remove all developer test rows before
launch. \\
\bottomrule
\end{tabularx}
\end{table}

\section{Questions for Advisor Review}

\begin{enumerate}[leftmargin=1.6em]
  \item Are the four research questions appropriately scoped, particularly
        RQ1's use of ``mental models'' and RQ4's current reliance on perceived
        auditability?
  \item Should an objective auditability scenario be added to strengthen RQ4,
        or should RQ4 remain focused on knowledge and self-reported confidence?
  \item Which independent data-quality and exclusion rules should be
        preregistered?
  \item What sample size is supported by the planned power analysis for the
        chain-by-placement interaction?
  \item Is the proposed Prolific recruitment of adult university students
        consistent with the intended population and IRB application?
  \item Is the combined task-and-survey burden acceptable after pilot timing,
        or should some open-text survey items be shortened or removed?
\end{enumerate}

\section{Summary}

The project now has a functioning behavioral prototype, a structured survey,
and a live data pipeline. Its central methodological strength is that it
separates two randomized causes---authentication-chain exposure and permission
placement---while holding the permission set constant. The behavioral task
captures what participants could see and what they did; the recall task and
survey capture what they remembered, believed, and understood. The main work
remaining is methodological finalization with the advisor and IRB, followed by
Qualtrics/Prolific integration, complete end-to-end validation, power analysis,
and a small pilot before recruitment.

\appendix
\section{Complete Follow-up Survey Instrument}

The following appendix presents the current 36-item Qualtrics instrument in the
order participants will receive it. Qualtrics will first capture
\texttt{pid}, \texttt{condition}, and \texttt{placement} as embedded data from
the task URL. The participant-code question is displayed only when
\texttt{pid} is missing.

\subsection*{Block 1: Entry and Eligibility}

\textbf{Q1. Participant code.}
Your participant code is captured automatically from the previous task. If the
field below is empty, please enter the code shown on the task completion
screen.

\textit{Response format:} Single-line text entry; hidden when
\texttt{pid} is populated.

\textbf{Q2. Age confirmation.}
I am 18 years of age or older.

\textit{Response options:} Yes; No.

\textit{Routing:} A No response ends the survey.

\textbf{Q3. Student status.}
I am currently enrolled as a student at a college or university.

\textit{Response options:} Yes; No.

\textit{Routing:} A No response sets \texttt{ineligible=true} and ends the
survey.

\subsection*{Block 2: Immediate Post-task Reflection}

\textbf{Q4. Realism.}
How realistic did the simulated account-access and permission task feel?

\textit{Response options:} (1) Not realistic at all; (2) Slightly realistic;
(3) Moderately realistic; (4) Very realistic; (5) Extremely realistic.

\textbf{Q5. Self-reported attention.}
During the task, how carefully did you read the permission screen before
making your decision?

\textit{Response options:} (1) I did not read it; (2) I skimmed it quickly;
(3) I read some of it; (4) I read most of it; (5) I read all of it carefully.

\textbf{Q6. Self-reported scrolling.}
Did you scroll through the full permission list during the task?

\textit{Response options:} Yes; No; I do not remember.

\textbf{Q7. Decision and broad-access reflection.}
Briefly explain why you chose to allow or cancel access during the task, and
whether anything about what the assistant asked for stood out to you.

\textit{Response format:} Open-text response.

\textbf{Q8. Suspicion check.}
Before the debrief, what did you think the study was mainly examining?

\textit{Response format:} Open-text response.

\subsection*{Block 3: Familiarity and Exposure}

\textbf{Q9. SSO frequency.}
In the past month, how often have you logged in to a site or app using an
existing account, such as ``Continue with Google,'' ``Continue with
Microsoft,'' or ``Continue with Apple''?

\textit{Response options:} (1) Never; (2) Rarely; (3) Monthly; (4) Weekly;
(5) Daily.

\textbf{Q10. Verification frequency.}
In the past month, how often have you had to verify your identity using a
second method, such as a code, phone prompt, or fingerprint?

\textit{Response options:} (1) Never; (2) Rarely; (3) Monthly; (4) Weekly;
(5) Daily.

\textbf{Q11. AI-assistant use.}
Have you used an AI assistant that can act on your behalf, such as reading
files, drafting or sending messages, scheduling events, or connecting to
another service?

\textit{Response options:} Yes, regularly; Yes, once or twice; No; Not sure.

\textbf{Q12. Technology comfort.}
How comfortable are you with technology in general?

\textit{Response options:} (1) Not at all comfortable; (2) Slightly
comfortable; (3) Moderately comfortable; (4) Very comfortable; (5) Extremely
comfortable.

\subsection*{Block 4: Security-behavior Baseline}

Questions 13--16 appear as one matrix with the shared response scale:
(1) Never, (2) Rarely, (3) Sometimes, (4) Often, and (5) Always.

\textbf{Q13. Link checking.}
When someone sends me a link, I open it only after checking where it actually
goes.

\textbf{Q14. Permission-review habit.}
I read what a website or app is asking for before I agree to it.

\textbf{Q15. Response to unusual behavior.}
I notice when a site or app is behaving in a way that seems off, and I stop
before continuing.

\textbf{Q16. Instructed-response check 1.}
To show you are reading carefully, please select ``Often'' for this item.

\textit{Scoring:} Pass when Often is selected.

\subsection*{Block 5: Mental Models Across Three Delegation Layers}

\textbf{Q17. SSO default access.}
When you sign in to an app using an existing account, such as Google,
Microsoft, or Apple, what do you assume the app receives by default?

\textit{Response options:}
\begin{itemize}[leftmargin=1.6em]
  \item Nothing; it only confirms I am a real person.
  \item Confirmation of my identity plus limited profile details the app asks
        for, such as name and email.
  \item Full access to that account, including messages and files.
  \item Not sure.
\end{itemize}

\textit{Keyed answer:} Confirmation of identity plus the limited profile
details requested.

\textbf{Q18. SSO and email scope.}
``Signing in to an app with my Google or Microsoft account automatically lets
that app read my email.''

\textit{Response options:} True; False; Not sure.

\textit{Keyed answer:} False.

\textbf{Q19. Open-text SSO understanding.}
In one or two sentences, what do you think happens when you sign in to an app
using an existing account like Google, Microsoft, or Apple?

\textit{Response format:} Open-text response, coded post hoc as accurate,
partially accurate, inaccurate, or off-topic.

\textbf{Q20. Purpose of MFA.}
What is the main thing two-step verification, such as a code or phone approval
after your password, is meant to protect against?

\textit{Response options:}
\begin{itemize}[leftmargin=1.6em]
  \item Weak passwords.
  \item Someone using a password that was stolen or guessed.
  \item Computer viruses.
  \item A slow or unreliable internet connection.
  \item Not sure.
\end{itemize}

\textit{Keyed answer:} Someone using a password that was stolen or guessed.

\textbf{Q21. SSO-replaces-MFA misconception.}
``If I use single sign-on, I do not really need two-step verification, because
the sign-on already covers it.''

\textit{Response options:} True; False; Not sure.

\textit{Keyed answer:} False.

\textbf{Q22. Persistence of AI access.}
When you allow an AI assistant to connect to one of your accounts, such as
email or calendar, what do you assume it can do afterward?

\textit{Response options:}
\begin{itemize}[leftmargin=1.6em]
  \item Only the single action I asked for, that one time.
  \item Whatever the granted access allows, until I remove that access.
  \item Only things it does while I am watching.
  \item Not sure.
\end{itemize}

\textit{Keyed answer:} Whatever the granted access allows, until the access is
removed.

\textbf{Q23. AI-assistant autonomy.}
``An AI assistant only does something at the exact moment I ask; it cannot act
on its own later.''

\textit{Response options:} True; False; Not sure.

\textit{Keyed answer:} False.

\textbf{Q24. Importance of permission clarity.}
How important is it to you that an AI assistant clearly explains what account
permissions it is requesting?

\textit{Response options:} (1) Not important at all; (2) Slightly important;
(3) Moderately important; (4) Very important; (5) Extremely important.

\subsection*{Block 6: Attention Across Authentication Steps}

\textbf{Q25. Attention across steps.}
When you go through several login or approval steps in a row, such as sign-in,
provider handoff, and MFA approval, how does your attention to each new screen
usually change?

\textit{Response options:}
\begin{itemize}[leftmargin=1.6em]
  \item I read each one just as carefully.
  \item I pay a little less attention to later ones.
  \item I pay a lot less attention to later ones.
  \item I usually approve or continue to get through them.
  \item Not sure.
\end{itemize}

\textbf{Q26. Login fatigue.}
How mentally tiring do repeated login and approval steps feel to you?

\textit{Response options:} (1) Not tiring at all; (2) Slightly tiring;
(3) Moderately tiring; (4) Very tiring; (5) Extremely tiring.

\subsection*{Block 7: Revocation and Auditability}

\textbf{Q27. Knowing where to revoke access.}
If you wanted to stop an app or AI assistant from being able to access your
account, would you know where to do that?

\textit{Response options:} Yes, confidently; I think so; No; I have never
thought about it.

\textbf{Q28. Revocation location.}
Briefly, where would you go to remove that access?

\textit{Response format:} Open-text response, coded as correct, partially
correct, or incorrect.

\textbf{Q29. Prior revocation experience.}
Have you ever reviewed or removed an app's or service's access to one of your
accounts?

\textit{Response options:} Yes; No; I am not sure what that means.

\textbf{Q30. AI-action attribution.}
If an AI assistant sent an email for you, do you think you could find a record
showing that the assistant sent it rather than you?

\textit{Response options:} Yes; No; Not sure.

\textbf{Q31. Instructed-response check 2.}
To show you are paying attention, please select ``False'' for this item:
``Most people sleep eight hours a night.''

\textit{Response options:} True; False; Not sure.

\textit{Scoring:} Pass when False is selected.

\subsection*{Block 8: Demographics and Closing}

\textbf{Q32. Major or field of study.}
What is your major or field of study?

\textit{Response format:} Open-text response, coded post hoc as technical,
non-technical, or mixed.

\textbf{Q33. Degree level.}
What is your current degree level?

\textit{Response options:} Undergraduate; Master's; Doctoral or PhD; Other;
Prefer not to say.

\textbf{Q34. Gender.} This question is optional.

\textit{Response options:} Woman; Man; Non-binary; Prefer to self-describe;
Prefer not to say.

\textbf{Q35. Prior cybersecurity coursework.}
Have you taken a course or formal training related to cybersecurity, privacy,
or computer security?

\textit{Response options:} Yes; No; Not sure.

\textbf{Q36. Final comment.}
Is there anything else you would like to tell us about the task, account
permissions, or AI assistants?

\textit{Response format:} Optional open-text response.

\subsection*{Closing Message}

Thank you for completing the survey. Your responses have been recorded. If you
have any questions about the study, please contact Kofi Ampomah at
\href{mailto:kampomah@crimson.ua.edu}{kampomah@crimson.ua.edu} or Kunal Sarna
at \href{mailto:ksarna@crimson.ua.edu}{ksarna@crimson.ua.edu}. For questions
about your rights as a research participant, contact the University of Alabama
Office for Research Compliance at 1-877-820-3066 or
\href{mailto:rscompliance@ua.edu}{rscompliance@ua.edu}.

\bibliographystyle{plain}
\bibliography{references}

\end{document}
